\section{Related work}
\label{sec:related}

Nothing in this library is a new algorithm. Every technique on the ladder has a
paper behind it, and the contribution here is the measurement discipline applied
to one part, not the invention. It is worth saying where each idea comes from,
because the interesting results below are the places where the received wisdom
does not survive contact with this GPU.

\paragraph{Blocked dense GEMM.} The layered blocking that every fast GEMM uses,
packing panels so that the innermost kernel streams from the fastest level of the
hierarchy, is Goto and van de Geijn's anatomy~\cite{goto2008}. The GPU restatement
of it, and the observation that register blocking rather than shared memory
tiling is what actually buys the throughput, is Volkov and Demmel~\cite{volkov2008}.
Nath, Tomov and Dongarra showed on Fermi that the winning tile is a property of
the shape and the part rather than a constant, and that autotuning over a family
of instantiations is how you find it~\cite{nath2010}; that argument is why Section
\ref{sec:mechanisms} ships six tile shapes rather than one.

\paragraph{Template libraries and compilers.} CUTLASS is the reference
implementation of the whole Ampere style multistage mainloop, and its CuTe layout
algebra is the vocabulary the swizzle and the fragment layouts in this project are
most easily described in~\cite{cutlass}. Section \ref{sec:cutlass} runs a CUTLASS
instantiation as a rung on my own ladder rather than citing it from a distance.
Triton attacks the same problem from the compiler side, generating the tiled
kernel from a higher level description~\cite{tillet2019}; I write the mainloop by
hand here because the point of the exercise is to see the machine, but Triton is
the honest answer to why most people should not.

\paragraph{Tensor cores and their numerics.} Markidis and colleagues measured what
the WMMA path actually costs numerically, and established the shape of the error
model this report reuses in Section \ref{sec:accuracy}: FP16 inputs with FP32
accumulate, where the input rounding term rather than the accumulation term
dominates~\cite{markidis2018}. Jia and colleagues dissected the tensor pipeline by
microbenchmarking, which is the method Section \ref{sec:sass} borrows at a smaller
scale when it reads the instruction mix out of the SASS instead of trusting the
source~\cite{jia2018}.

\paragraph{Work decomposition.} Stream-K is Osama, Merrill, Cecka, Garland and
Owens~\cite{osama2023}: instead of cutting every tile the same way, cut only the
work that would have left SMs idle, which is exactly the wave quantization problem
Section \ref{sec:gap} measures at 1024 cubed. The implementation in
\texttt{src/gemm} follows that paper's decomposition, with the partial reduction
protocol described in Section \ref{sec:mechanisms}.

\paragraph{Fusion.} FlashAttention is the clearest statement of the fusion
argument this report uses for its epilogue study: when a pass is memory bound, the
win comes from not making the pass at all, and the way to see that is to count
bytes rather than FLOP~\cite{dao2022}. The bias plus ReLU epilogue in Section
\ref{sec:roofline} is the same argument at a much smaller scale.

\paragraph{Sparse.} Bell and Garland set the CSR baselines that the naive and warp
per row kernels in Section \ref{sec:sparse} are~\cite{bell2009}. Merrill and
Garland's merge based formulation removes the load imbalance those two still
carry, by giving every thread an equal share of the merge of row offsets with
nonzero indices~\cite{merrill2016spmv}. Kreutzer and colleagues' SELL-C-sigma is
the other direction, changing the storage format so the inner loop coalesces at
the cost of padding~\cite{kreutzer2014}. Both are implemented here.

\paragraph{Scan.} Blelloch's work efficient formulation is the classical
result~\cite{blelloch1990}, and it is a rung on the scan ladder rather than the top
of it. The top rung is Merrill and Garland's decoupled look-back, which reaches the
memory bound of the problem by making one pass instead of three and paying for it
with a device scope handshake between tiles~\cite{merrill2016scan}.

\paragraph{FFT.} The Stockham autosort formulation, and the four step algorithm
that makes a transform larger than shared memory into a sequence of resident sub
transforms, are both in Van Loan's framework treatment~\cite{vanloan1992}. Section
\ref{sec:fft} follows that formulation directly.

\paragraph{This part specifically.} Two sources bear on why a hand written kernel
has room against cuBLAS here at all. The RTX Blackwell whitepaper gives the
machine~\cite{blackwell}, and an independent analysis of cuBLAS on consumer
Blackwell reports it dispatching Ampere generation
\texttt{cutlass\_80\_tensorop\_h16816gemm} kernels for FP16 GEMM~\cite{cloudrift}.
I have not reproduced that dispatch analysis on my own card, so I carry it as
context for why the gap exists rather than as a result of mine; Section
\ref{sec:limitations} says so again.
